\section{Clocked traces, invariants, and termination}\label{sec:traces}

\subsection{The trace normal form}

Many finite mathematical proofs can be normalized to a primitive-recursive state machine:
\[
 s_0=\mathsf{init}(\vec x),\qquad
 s_{t+1}=\mathsf{step}(t,s_t,\vec x).
\]
The state is one natural code, typically a nested pairing of fields.  A public theorem reads an
observation from \(s_T\) at an explicit horizon \(T(\vec x)\).  The proof burden is:

\begin{enumerate}
\item defining equations for the trace;
\item a quantifier-free invariant \(I(t,s_t,\vec x)\);
\item a base derivation and a preservation derivation;
\item a finalizer deriving the desired formula from \(I(T,s_T,\vec x)\).
\end{enumerate}

All rich mathematical structure is represented by finite fields and certificate predicates inside
\(I\).  Quantifiers needed by the informal proof are either bounded-reflected before induction or
moved into the finalizer.

\subsection{Generic invariant theorem}

Let \(r(t,\vec x)\) be defined by primitive recursion from a base term \(b(\vec x)\) and step term
\(s(t,r,\vec x)\).  Let \(I(t,\vec x)\) abbreviate a quantifier-free formula in which occurrences of
the current state are replaced by \(r(t,\vec x)\).

\begin{theorem}[Clocked trace reflection]\label{thm:trace}
If
\[
 []\vdash_{\PRAAx} I(0,\vec x)
\]
and
\[
 I(t,\vec x)\vdash_{\PRAAx} I(t+1,\vec x),
\]
then
\[
 []\vdash_{\PRAAx} I(t,\vec x)
\]
for all \(t\), with the other variables treated as parameters.  Hence every explicit horizon
\(T(\vec x)\) may be substituted into the universal result.
\end{theorem}

\begin{proof}
Apply quantifier-free induction to the formula \(I\).  The base and successor premises are exactly
those required by \code{PRA.qfInd} (or by the open-variable wrapper \code{IPRA.ind}).  Universal
elimination specializes the result to \(T(\vec x)\).  The defining equations of \(r\), proved by
\code{rec_zero} and \code{rec_succ}, are used in the base and step derivations to rewrite the state.
\end{proof}

\subsection{State design}

A robust state layout contains:

\begin{itemize}
\item the mathematical accumulator;
\item the loop counter when it is not already the recursion variable;
\item a failure bit that becomes one permanently after a local check fails;
\item resource parameters whose values would otherwise change the induction formula;
\item optional certificate fragments or a rolling hash/summary when the final theorem requires a
      trace certificate.
\end{itemize}

A failure accumulator is particularly useful.  Let \(e_t\in\{0,1\}\) indicate whether the local
condition at stage \(t\) succeeds, and define
\[
 F_0=0,\qquad F_{t+1}=\operatorname{sg}(F_t+(1\mathbin{\dot-}e_t)).
\]
Then \(F_T=0\) is equivalent to every local condition succeeding for \(t<T\).  The invariant remains
quantifier free, while bounded universal meaning is recovered by a reflection theorem.

\subsection{Measure-decreasing termination}

The pinned repository already proves an object theorem for stopped runs.  Given primitive-recursive
\(\mathsf{done}\), \(\mathsf{step}\), and measure \(m\), define a run that remains fixed after
termination.  Under the hypothesis
\[
 \forall s\,[\mathsf{done}(s)=0\to m(\mathsf{step}(s))<m(s)],
\]
the checked theorem proves that the run is terminal by time \(m(s_0)+1\).

\begin{corollary}[Fuel extraction]
A deterministic local algorithm whose nonterminal steps strictly decrease an explicit natural
measure can be compiled to a total fixed-horizon trace without postulating general well-founded
recursion.
\end{corollary}

\begin{proof}
Use the checked termination theorem with horizon \(m(s_0)+1\).  The trace is primitive recursive
because it is a stopped bounded iteration.  Once terminality is known, a finalizer reads the result.
\end{proof}

Lexicographic or multicomponent measures can be encoded into one natural by a proved ranking
function, or handled by an explicit finite capacity bound.  The measure theorem is a producer if its
decrease proof uses domain-specific mathematics; the generic termination step is already checked.

\subsection{Changing precision}

Suppose step \(t\) needs precision \(p_t\), and an informal proof chooses \(p_{t+1}\) after seeing the
current state.  Directly writing \(I_{p_t}(t)\) would change the induction formula.  There are two
correct transformations.

\paragraph{Reverse demand pass.}
For a fixed final query, compute all required precisions before object induction.  Put the resulting
finite table into a code and let the state read its \(t\)-th entry.

\paragraph{Precision as state.}
Add \(p_t\) to the state, define its primitive-recursive update, and strengthen the invariant to
assert that the current mathematical data meet the current precision.  The formula is fixed because
``current precision'' is a term projection, not a meta-level index.

Both transformations turn a family of informal invariants into one quantifier-free formula.

\subsection{Proof-size control}

Unrolling \(T\) local steps produces a proof term linear or worse in \(T\).  The trace theorem uses
one induction derivation independent of the numeric horizon.  The final proof contains the fixed
base and step proofs plus a term denoting \(T\).  For large equational calculations, a native-clause
backend can share common fragments, but the main asymptotic improvement already comes from primitive
recursion plus one induction.

\subsection{Nondeterministic traces}

A relational computation is represented in one of three ways:

\begin{enumerate}
\item a sealed whole-trace certificate checked by a primitive-recursive verifier;
\item a backward viability table followed by deterministic least choices;
\item a cloud-valued state whose public observer is stable under all legal successors.
\end{enumerate}

The first is simplest when a producer supplies a complete path.  The second exposes a canonical path
and therefore requires finite bounds.  The third is strongest constructively when only an invariant
observation matters.

\subsection{Lean integration pattern}

A theorem-specific trace module should contain, in dependency order:

\begin{enumerate}
\item the state projections and substitution lemmas;
\item the base and step \(\PTm\) definitions and their \code{rec_zero}/\code{rec_succ} equations;
\item the quantifier-free invariant and a proof of \code{Fm.QF};
\item base and preservation derivations;
\item one call to \code{PRA.qfInd} or \code{IPRA.ind};
\item a finalizer plan and schedule node.
\end{enumerate}

This layout localizes the only induction proof and makes residual producers visible in the context of
base, step, or finalizer nodes.
